\documentclass{article}

% --- load packages ---
\usepackage[margin=1in]{geometry}
\usepackage{amsmath}
\usepackage{amssymb}
\usepackage[colorlinks,bookmarks,bookmarksnumbered,allcolors=blue]{hyperref}
\usepackage{graphicx}
\usepackage{booktabs}
\usepackage[capitalise]{cleveref}
\usepackage[section]{placeins}
\usepackage{cite}
\usepackage{doi}
\usepackage{listings}
\usepackage{xcolor}

\lstset{
  basicstyle=\ttfamily\small,
  breaklines=true,
  frame=single,
  backgroundcolor=\color{gray!8}
}

\begin{document}

\title{Project 1: Proof-Producing CDCL SAT Solver}
\author{ZHANG Boyuan\\
  \small CSC\_54656 EP --- Satisfiability and SMT Solving
}
\date{}
\maketitle

%-----------------------------------------------------------------------
\section*{Abstract}
%-----------------------------------------------------------------------
We extend the CDCL solver developed in Tutorial~2 to emit a
\emph{resolution proof} for every unsatisfiable instance, and implement
an independent linear-time \emph{proof checker} that validates the proof
against the original formula.
The design follows the na\"ive strategy of~\cite{biere2008picosat}:
every resolution step performed during conflict analysis is recorded
verbatim, with no preprocessing steps and no clause deletion tracking.
This keeps the proof format minimal, the checker trivial to implement,
and the correctness argument transparent.
Experiments on 12 instances---including the pigeon-hole benchmark
family up to $n=8$---confirm that the solver and checker are correct.

%-----------------------------------------------------------------------
\section{Problem Description}
%-----------------------------------------------------------------------

\subsection{From Tutorial 2 to Project 1}

Tutorial~2 asked us to implement a complete CDCL SAT solver with
Boolean Constraint Propagation (BCP), 1st-UIP conflict analysis,
non-chronological backtracking, and clause learning.
The solver can answer \textsc{Sat} (returning a satisfying
assignment) or \textsc{Unsat}.

Project~1 extends this with a \emph{proof of unsatisfiability}.
This matters for two reasons.
First, from a \emph{trust} perspective: when a solver outputs
\textsc{Unsat}, no assignment exists as a witness---the only way to
gain confidence is to accompany the answer with a proof that a
trusted third party can check.
Second, from a \emph{correctness} perspective: bugs in SAT solvers are
known to produce wrong \textsc{Unsat} answers~\cite{biere2008picosat};
a checkable proof makes such errors detectable.

\subsection{Resolution as a Proof System}

Resolution is the canonical proof system for propositional
unsatisfiability (\cref{eq:resolution}).
Given two clauses $(C_1 \vee \neg p)$ and $(C_2 \vee p)$ that contain
complementary occurrences of literal $p$ (the \emph{pivot}), their
\emph{resolvent} is $C_1 \vee C_2$, which is a logical consequence of
both parent clauses.

\begin{equation}
  \frac{C_1 \vee \neg p \quad C_2 \vee p}{C_1 \vee C_2}
  \label{eq:resolution}
\end{equation}

A \emph{resolution refutation} is a finite sequence of such steps
starting from the input clauses and deriving the empty clause
$\bot$, which has no satisfying assignment---thereby witnessing
unsatisfiability.

The connection to CDCL is direct: every \emph{conflict analysis} step
already computes a resolution (resolving the current conflict clause
against a reason clause from the trail to eliminate one literal).
Logging each such step is therefore a natural and low-overhead extension.

%-----------------------------------------------------------------------
\section{Implementation}
%-----------------------------------------------------------------------

\subsection{Proof Format Design}

\paragraph{Design choice.}
We chose the simplest possible line-oriented plain-text format.
Each line encodes exactly one resolution step:

\begin{lstlisting}
p proof <n>
<new_id> <pivot> <id1> <id2> : <lit1> <lit2> ... 0
\end{lstlisting}

\noindent
The header \texttt{p proof n} declares that initial clauses carry IDs
$1\ldots n$ (in DIMACS order).
Each subsequent line gives the new clause~ID, the pivot literal,
the two parent IDs, and the resolvent literals terminated by zero.
The empty clause is represented by a bare \texttt{0}.

\paragraph{Trade-offs.}
The format is deliberately \emph{redundant}: the resolvent is fully
spelled out even though a checker could recompute it.
This makes verification a pure equality check---no resolution
computation is needed---and makes the proof human-readable for debugging.
The cost is a larger file size compared to a compact binary format such
as DRAT~\cite{biere2008picosat}, but for the problem sizes considered
(up to 4.6\,MB for hole8) this is acceptable.
A single forward pass over the file with no backtracking suffices for
both production and checking.

\subsection{Equipping the Solver with Proof Production}

Three components were added to the Tutorial~2 CDCL skeleton in
\texttt{cdcl.py}.

\paragraph{Proof logger (trade-off: branch-free design).}
\texttt{ProofLogger} wraps a file handle and exposes one method:
\texttt{log\_step(new\_id, pivot, id1, id2, resolvent)}.
A companion \texttt{\_NullLogger} implements the identical interface
as a no-op.
The key design decision is that \texttt{CDCL} always holds a logger
object (either real or null); the main solver loop contains
\emph{no conditionals for proof mode}.
This avoids performance regressions when proof production is disabled
and eliminates a class of bugs where a conditional branch is
accidentally omitted.
Proof production is activated by passing \texttt{proof\_file} to the
constructor; the default \texttt{None} preserves full backward
compatibility with the Tutorial~2 interface.

\paragraph{Clause ID tables (trade-off: dual lookup strategy).}
\texttt{ClauseDb} maintains two dictionaries:
\begin{itemize}
  \item \texttt{clause\_id[\,id(list)]} maps \emph{Python object
    identity} to proof ID.
    This is the fast path: \texttt{deduce()} returns the exact list
    object stored in \texttt{ClauseDb.clauses}, so the identity lookup
    is an $O(1)$ dictionary hit.
  \item \texttt{content\_to\_id[\,tuple(sorted(lits))]} provides a
    \emph{content-based fallback}.
    When a reason clause is accessed via a different code path (e.g.\
    after \texttt{learn()} appends a new list), only its content is
    guaranteed to match.
\end{itemize}
Initial clauses receive IDs $1\ldots n$; learned clauses are assigned
IDs from $n{+}1$ upward via a shared counter.
This dual-map strategy avoids a full scan of all clauses to find a
proof ID, keeping each lookup to $O(1)$.

\paragraph{Proof derivation and the level-0 edge case.}
\texttt{analyze()} logs one resolution step for each literal eliminated
during 1st-UIP derivation, via the helper \texttt{\_proof\_resolve\_step()}.
This produces a learned clause whose proof ID is registered before the
clause is passed to \texttt{learn()}.

A subtle edge case arises when the solver must return \textsc{Unsat}
but the derived learned clause is not yet the empty clause: after
backtracking to level~0 and running BCP again, a new conflict is
detected whose resolution to $\bot$ requires resolving against BCP-forced
unit clauses on the trail.
The solution is to \emph{postpone} empty-clause derivation: instead
of calling \texttt{\_finish\_proof\_to\_empty} inside \texttt{analyze},
\texttt{solve()} detects any conflict at decision level~0 and calls
\texttt{\_finish\_proof\_to\_empty} at that point.
At level~0 the trail contains only BCP-forced entries (no decisions),
so every literal in the working clause has a reason clause available,
and the loop terminates with the empty clause.

\subsection{Proof Checker (\texttt{proof\_checker.py})}

\paragraph{Independence.}
The checker shares no code with \texttt{cdcl.py}.
It can be used to validate proofs from \emph{any} solver that emits
the same format, not just ours.

\paragraph{Algorithm.}
The checker reads the DIMACS instance into a dictionary
$\textit{clause\_db}\colon \mathbb{Z} \to \mathcal{P}(\mathbb{Z})$
(mapping clause ID to literal set), then makes one forward pass:

\begin{enumerate}
  \item Look up $C_1 = \textit{clause\_db}[\textit{id1}]$ and
        $C_2 = \textit{clause\_db}[\textit{id2}]$ in $O(1)$.
  \item Verify $-\textit{pivot} \in C_1$ and $+\textit{pivot} \in C_2$.
  \item Compute the expected resolvent
        $(C_1 \setminus \{-p\}) \cup (C_2 \setminus \{p\})$
        and compare with the declared literals.
  \item Store the validated resolvent under \textit{new\_id}.
\end{enumerate}

After the scan, the checker asserts that the last produced clause is
the empty set.
Time complexity is $O(N \cdot K)$ where $N$ is the number of steps
and $K$ the maximum clause size: linear in the total proof size.
The checker never backtracks and never re-reads a line.

%-----------------------------------------------------------------------
\section{Experimental Evaluation}
%-----------------------------------------------------------------------

\subsection{Test Suite Design}

The test suite covers three categories, all runnable via
\texttt{run\_experiments.py}:

\begin{itemize}
  \item \textbf{SAT instances} (\texttt{prop.cnf}, \texttt{simple.cnf},
    \texttt{single\_literal.cnf}, \texttt{empty\_lines.cnf},
    \texttt{with\_comments.cnf}, \texttt{multiline.cnf},
    \texttt{large.cnf}): verify that the DIMACS reader handles varied
    formatting (blank lines, inline comments, multi-line clauses) and
    that satisfiable instances are not mistakenly declared unsatisfiable.
  \item \textbf{UNSAT edge cases}:
    \texttt{border\_level0.cnf} ($\{x_1\} \wedge \{\neg x_1\}$)
    is the minimal UNSAT instance, testing the level-0 conflict path
    where no decision is ever made.
    \texttt{DPLL example2} (2 variables, 4 clauses) tests the minimal
    non-trivial proof of 4 resolution steps.
  \item \textbf{Pigeon-hole benchmarks} $\text{PHP}^{n+1}_n$:
    exponentially hard for resolution~\cite{cook1976}, used to stress-test
    proof size and checker performance.
\end{itemize}

\subsection{Experiment Runner (\texttt{run\_experiments.py})}

Running experiments manually is error-prone and not reproducible.
\texttt{run\_experiments.py} automates the full pipeline:
(1) invoke the solver on each instance,
(2) write the proof file for \textsc{Unsat} instances,
(3) call \texttt{check\_proof()} directly (no subprocess overhead) to
validate the proof,
(4) print a formatted results table.
Three modes are provided:

\begin{lstlisting}
python run_experiments.py                 # SAT + UNSAT edge cases
python run_experiments.py --pigeon        # add hole6/7/8
python run_experiments.py --hole 6 7 8    # choose specific instances
python run_experiments.py --check-only    # verify existing proof files
python run_experiments.py --clean         # delete generated proof files
\end{lstlisting}

\subsection{Results}

\Cref{tab:results} summarises results across all 12 test instances
(12/12 correct).

\begin{table}[htb]
\centering
\caption{Full results (12 instances, 12/12 correct).
  Proof information is reported only for \textsc{Unsat} instances.}
\label{tab:results}
\begin{tabular}{@{}lrrrrrl@{}}
\toprule
Instance & Vars & Clauses & Result & Solver (s) & Proof steps & Checked \\
\midrule
\texttt{prop.cnf}            &  3 &    3 & SAT   & $<$0.001 & ---      & ---   \\
\texttt{simple.cnf}          &  3 &    2 & SAT   & $<$0.001 & ---      & ---   \\
\texttt{single\_literal.cnf} &  5 &    4 & SAT   & $<$0.001 & ---      & ---   \\
\texttt{empty\_lines.cnf}    &  2 &    2 & SAT   & $<$0.001 & ---      & ---   \\
\texttt{with\_comments.cnf}  &  4 &    3 & SAT   & $<$0.001 & ---      & ---   \\
\texttt{multiline.cnf}       &  6 &    3 & SAT   & $<$0.001 & ---      & ---   \\
\texttt{large.cnf}           & 10 &    5 & SAT   & $<$0.001 & ---      & ---   \\
\texttt{border\_level0.cnf}  &  1 &    2 & UNSAT & $<$0.001 & 1        & VALID \\
\texttt{DPLL example2}       &  2 &    4 & UNSAT & $<$0.001 & 4        & VALID \\
\texttt{hole6.cnf}           & 42 &  133 & UNSAT & 0.27     & 5\,124   & VALID \\
\texttt{hole7.cnf}           & 56 &  204 & UNSAT & 2.66     & 18\,206  & VALID \\
\texttt{hole8.cnf}           & 72 &  297 & UNSAT & 26.2     & 61\,077  & VALID \\
\bottomrule
\end{tabular}
\end{table}

\subsection{Discussion}

All 12 instances produce the correct result (7 SAT, 5 UNSAT), and all
5 proofs are confirmed valid by the independent checker.

The \texttt{border\_level0.cnf} instance ($\{x_1\} \wedge \{\neg x_1\}$)
produces a single-step proof: clause~1 and clause~2 are resolved on
pivot~$x_1$ directly to obtain the empty clause.
This confirms that the level-0 edge case is handled correctly without
invoking any decision or backtracking.

The pigeon-hole results confirm the known exponential lower bound on
resolution proof size~\cite{cook1976}: proof steps grow from
$5\,124$ (hole6) to $18\,206$ (hole7) to $61\,077$ (hole8), a factor
of roughly $3\times$ per unit increase in $n$.
Despite the na\"ive implementation (no watched literals, no restart,
no VSIDS heuristic), all instances are solved correctly.
The checker verifies the 4.6\,MB hole8 proof in under 0.25\,s,
confirming the linear-time complexity of the checking algorithm.

The key design decision---keeping the proof format minimal
(one line per step, redundant resolvent spelled out)---pays off:
the checker's correctness argument is a simple loop invariant with
no auxiliary data structures beyond the clause dictionary.

%-----------------------------------------------------------------------
\section{Conclusion}
%-----------------------------------------------------------------------

We extended a Tutorial~2 CDCL solver with proof production and
implemented an independent linear-time checker.
The branch-free logger design, dual-map clause ID lookup, and
postponed empty-clause derivation are the three main technical
contributions.
All 12 test instances are handled correctly, including the level-0
conflict edge case and pigeon-hole instances up to $n=8$
(61\,077 proof steps, checker time $<$0.25\,s).

Future work could adopt the DRAT format~\cite{biere2008picosat}, which
supports clause deletion and non-resolution inferences used by
competition-grade solvers, at the cost of a more complex checker.

%-----------------------------------------------------------------------
\bibliographystyle{unsrt}
\begin{thebibliography}{9}

\bibitem{biere2008picosat}
A.~Biere.
\newblock {PicoSAT} essentials.
\newblock \emph{Journal on Satisfiability, Boolean Modeling and Computation},
  4:75--97, 2008.

\bibitem{cook1976}
S.~Cook and R.~Reckhow.
\newblock The relative efficiency of propositional proof systems.
\newblock \emph{Journal of Symbolic Logic}, 44(1):36--50, 1979.

\bibitem{handbook}
A.~Biere, M.~Heule, H.~van~Maaren, and T.~Walsh, editors.
\newblock \emph{Handbook of Satisfiability}.
\newblock IOS Press, 2nd edition, 2021.

\end{thebibliography}

\end{document}
